% !TEX root = main.tex
\begin{table}[!t]
\centering
\resizebox{\linewidth}{!}{
\begin{threeparttable}
\caption{Comparing Hash-Based Multilinear Polynomial Commitment Schemes}
\label{table:pcs}
\begin{tabular}
	{>{\arraybackslash}p{2.8cm}
	>{\centering\arraybackslash}p{3.3cm}
	>{\centering\arraybackslash}p{4.4cm}
	>{\centering\arraybackslash}p{4.0cm}
	}
	\hline
	 & \makecell{Commitment time}
	 & \makecell{Opening time}
	 & \makecell{Proof size}
	\\
	\hline
	Brakedown\cite{crypto/GolovnevLSTW23}
	 & $25.5n\Field,\, 2n\hashes$
	 & $O(n)\Field$
	 & $O(\lambda n)\Field$
	\\
    \hline
	BaseFold\cite{crypto/ZeilbergerCF24}
	 & $O(n\log n)\Field,\, 8n \log_2 n\hashes$
	 & $O(n)\Field$
	 &  $\log_2^2 n \cdot \lambda  \mpath$
	\\
	\hline
	Blaze(RAA)~\cite{eurocrypt/BrehmCFRRZ25}
	 & $8n\Field_{+},\, 4n\hashes$
	 & $O(n)\Field,\, O(n)\hashes$
	 & $(11+ \log_2 n) \cdot \lambda\mpath$
	\\
	\hline
	FRI
	 & $4 n\log n \Field,\, 4 n \hashes$
	 & $O(n)\Field,\, O(n)\hashes$
	 & $O(\log n \cdot \lambda) \mpath$
	\\
	\hline
    WHIR~\cite{eurocrypt/ArnonCFY25}
	 & $2 n\log n \Field,\, 4n \hashes$
	 & $\tilde{O}(n)\Field,\, O(n)\hashes$
	 & $O(\log \log n \cdot \lambda) \mpath$
     \\
	\hline
    TensorSwitch(RAA)
	 & $42n\Field_{+},\, 3n \Field,\, 16n\hashes$
	 & $6n\Field, (\lambda n)^{0.51} \hashes$
	 & $19\log_2n \cdot \lambda \hashes$
	\\
	\hline
    \hline
    \methodname(RAA)
	 & $8n\Field_{+}, n\log_2n\Field, \, 4n\hashes$
	 & $(8n/log_2n)\Field_{+}, n\Field,\, (4n/log_2n)\hashes$
	 & $(16 + O(\log\log n)) \cdot \lambda \mpath$
	\\
    \hline
\end{tabular}
\begin{tablenotes}
	\item[Note] 
	$\Field$ denote field elements or operations, $\hashes$ are pairing and hashing.
	$\epsilon \ge 2$ is a constant integer.
    For Basefold, we report the result for a code with $1/8$ rate.
    Brakedown is instantiated with $\rho = 0.58$.
    The RAA code in Blaze, TensorSwitch(Blaze), TensorSwitch+(Blaze), and DualSwitch(Blaze) is instantiated with rate $\rho = 1/4$.
    For WHIR and \tensorswitchplus and \methodname,
    we assume using reed-solomon code of rate $1/4$.
\end{tablenotes}
\end{threeparttable}
}
\end{table}
